% ============================================================
% 00_packages.tex — Pakete und Paket-Level-Setup
% ============================================================

\usepackage[T1]{fontenc}
\usepackage[utf8]{inputenc}
\usepackage[sfdefault,lf]{carlito} % Sans-Serif für Fliess-Text (bessere Lesbarkeit), Mathe bleibt Serif

\usepackage[landscape, a4paper, margin=1.5mm]{geometry}
\usepackage{microtype}
\usepackage{multicol}

\usepackage{mathtools}             % Superset von amsmath: DeclarePairedDelimiter, coloneqq etc.
\let\Bbbk\relax                    % Verhindert Konflikt zwischen newpxmath und amssymb
\usepackage{amssymb}
\usepackage{siunitx}
\sisetup{
  mode=match,                      % Einheiten passen sich der Schriftart an (wichtig mit carlito)
  propagate-math-font=true,
  reset-text-family=false,
  exponent-product = \cdot,
  output-decimal-marker = {,},
  per-mode = symbol
}

\usepackage[table]{xcolor}         % [table] aktiviert \rowcolor, \columncolor in Tabellen
\usepackage{array}                 % Erweiterte Spaltentypen und \arrayrulecolor
\usepackage{tabularx}              % Tabellen mit \linewidth-Fit und X-Spalten
\usepackage{tabularray}            % Kompakte valuegrid-Tabellen
\usepackage{xstring}
\usepackage{ragged2e}              % Silbentrennung bei Flattersatz in schmalen Spalten
\usepackage{enumitem}              % Kompakte Listen via \setlist

\usepackage{graphicx}
\usepackage{tcolorbox}
\tcbuselibrary{skins, breakable, raster}

\usepackage{makeidx}               % Stichwortverzeichnis (Setup in 65_index)
\usepackage{hyperref}
